\section{Теория на корутините и асинхронното изпълнение в C++}
\label{sec:coroutines_async_theory}

Корутините са езиков механизъм за описване на прекъсваеми изчисления, при които изпълнението може да бъде спряно и по-късно възобновено без загуба на локалното състояние. В стандарта C++20 корутините влизат в основния език, а не се реализират като външна библиотечна симулация~\cite{cpp20standard,p0912r5,p1063r0}. Това е важно за настоящата работа, защото позволява асинхронният слой на библиотеката да се изгражда върху compile-time дефинирани правила за типове, lifetime и control flow, вместо върху ad hoc callback структура.

В тази глава корутините се разглеждат не като удобен синтаксис, а като изпълнителен модел, който стои между чисто синхронното блокиращо изпълнение и тежките нишки на операционната система. Целта не е да се покаже само как се използват \code{co_await} и \code{co_return}, а да се изведе защо компилаторното преобразуване към state machine прави този подход особено подходящ за библиотека, която и без това се основава на статични типове, trait специализации и compile-time валидиране.

Този теоретичен пласт е необходим и поради още една причина. В контекста на \ORMName{} корутините не са изолиран езиков експеримент, а основа за единен асинхронен интерфейс над backend-и с много различни operational характеристики. За да бъде такава архитектура научно защитима, трябва ясно да се разграничат езиковият механизъм, runtime инфраструктурата и конкретните драйверни execution semantics.

\subsection{Процеси, нишки, задачи и корутини}

Процесът е защитен адресен контекст на операционната система, нишката е kernel-планирана последователност от инструкции, а задачата е по-общо програмно понятие за единица работа. Корутината не е нито процес, нито нишка. Тя е кооперативна единица за изпълнение, чието спиране и подновяване се управлява от програмния модел и generated runtime структура, а не директно от scheduler-а на операционната система.

Това разграничение е съществено, защото библиотеката комбинира корутини и \code{ThreadPool}, без да приравнява двете абстракции. Нишката е скъпа kernel ресурсна единица, свързана с отделен stack, scheduling и синхронизационни примитиви. Корутината е значително по-лека логическа единица, която пази само състоянието, необходимо за спиране и продължаване на конкретно изчисление. Следователно преминаването към coroutine-базиран модел не означава „повече нишки“, а по-структуриран начин за описване на това къде едно изчисление може да изчака външно събитие.

\begin{table}[H]
\centering
\caption{Съпоставка между основните изпълнителни единици}
\label{tab:execution_units_comparison}
\begin{tabularx}{\textwidth}{L{2.8cm}L{4.0cm}Y}
\toprule
\textbf{Единица} & \textbf{Същност} & \textbf{Значение за библиотеката} \\
\midrule
Процес & Изолиран адресен контекст с ресурси на операционната система & Не е пряка единица на ORM архитектурата, но определя границата на runtime средата \\
Нишка & Kernel-планиран поток на изпълнение със собствен stack & Използва се от \code{ThreadPool} за offload на блокиращи операции \\
Корутина & Кооперативно прекъсваемо изчисление със запазено локално състояние & Осигурява формалния модел за \code{Task<T>} и за асинхронната композиция \\
\bottomrule
\end{tabularx}
\end{table}

Таблица~\ref{tab:execution_units_comparison} подчертава, че корутината е логическа, а не kernel единица. Именно това позволява в една и съща нишка да се редуват множество suspend/resume последователности, без всяка от тях да изисква отделен stack и отделна нишкова идентичност.

\subsection{Синтаксис, awaitable protocol и езикови binding-и}

Ключовите езикови оператори \code{co_await}, \code{co_return} и \code{co_yield} маркират, че дадена функция следва coroutine semantics. От гледна точка на архитектурата най-важен е \code{co_await}, защото именно той описва точката, в която текущото изчисление може да отстъпи изпълнението, докато външна операция --- например база данни или scheduling към thread pool --- не стане готова. Така синтаксисът на потребителя остава близък до линейния код, докато изпълнителният модел става асинхронен.

За текущата дипломна работа е важно, че тази езикова форма не е динамично „откривана“ в runtime, а е компилаторно вързана към return type-а на функцията и към \code{promise\_type}, който \code{std::coroutine\_traits} извежда за конкретния подпис~\cite{cpp20standard,p0912r5}. Това означава, че начинът, по който една корутина създава резултат, предава грешка или съхранява continuation, може да бъде формализиран още на ниво типова система.

На практика \code{co_await} работи чрез awaitable protocol. Обектът, който се await-ва, участва в тристъпкова схема: \code{await\_ready()} определя дали спиране изобщо е необходимо, \code{await\_suspend()} получава coroutine handle и решава как да организира възобновяването, а \code{await\_resume()} връща резултата или повдига грешка. Точно тази тройка прави възможно различни източници на асинхронност --- thread pool задачи, event loop readiness, callback-и от драйвер --- да бъдат уеднаквени под общ \code{co_await} синтаксис.

\subsection{Корутината като state machine}

Компилаторът не „изпълнява магия“, а преобразува coroutine function в state machine с ясно определени suspend points, resume transitions и frame storage~\cite{p0912r5,p1063r0}. Това преобразуване е от особено значение за настоящата теза, защото показва, че compile-time подходът на библиотеката не се изчерпва само с моделиране на заявки. Същият езиков инструментариум позволява и асинхронният control flow да бъде формализиран чрез типове и статични правила, вместо да се описва чрез неструктурирани callback вериги.

Концептуално една coroutine функция преминава през няколко ясно разграничени фази: създаване на frame и \code{promise\_object}, начално изпълнение, достигане до suspend point, външно организирано възобновяване и финално завършване. Това е съществено, защото correctness вече не зависи само от „какво прави кодът“, а и от това кога е допустимо той да бъде продължен и кой носи отговорност за продължаването.

\begin{figure}[H]
\centering
\begin{tikzpicture}[node distance=1.5cm and 1.8cm]
    \node[ormblock, text width=3.1cm, align=center] (call) {Извикване на coroutine функция};
    \node[ormblock, text width=3.4cm, align=center, right=of call] (frame) {Създаване на coroutine frame и \code{promise\_object}};
    \node[ormblock, text width=2.8cm, align=center, right=of frame] (run) {Начално изпълнение на тялото};

    \node[ormaccent, text width=2.9cm, align=center, below=of frame] (await) {\code{co\_await} върху awaitable};
    \node[ormblock, text width=3.0cm, align=center, right=of await] (suspend) {Състояние на спряна корутина};
    \node[ormblock, text width=3.2cm, align=center, below=of suspend] (source) {Източник на resume:\\thread pool / event loop / драйверен callback};
    \node[ormaccent, text width=2.8cm, align=center, left=of source] (resume) {Възобновяване чрез coroutine handle};

    \node[ormblock, text width=2.9cm, align=center, left=of resume] (finish) {Продължаване на тялото и \code{co\_return}};
    \node[ormblock, text width=3.0cm, align=center, below=of finish] (final) {\code{final\_suspend}, continuation, резултат/грешка};

    \coordinate (run_suspend_mid) at ($(run)!0.5!(suspend)$);

    \draw[ormline] (call) -- (frame);
    \draw[ormline] (frame) -- (run);
    \draw[ormline] (run) -- (run |- run_suspend_mid) -- (await |- run_suspend_mid) -- (await.north);
    \draw[ormline] (await) -- (suspend);
    \draw[ormline] (suspend) -- (source);
    \draw[ormline] (source) -- (resume);
    \draw[ormline] (resume) -- (finish);
    \draw[ormline] (finish) -- (final);
\end{tikzpicture}
\caption{Концептуален lifecycle на корутината като frame, suspend points и външно управлявано възобновяване}
\label{fig:coroutine_lifecycle}
\end{figure}

Фигура~\ref{fig:coroutine_lifecycle} показва защо корутинният модел е естествен за библиотека с ясно отделени слоеве. Логическият код на потребителя остава линеен, но реалното управление на изпълнението преминава през формални преходи, които могат да бъдат свързани с конкретни runtime механизми.

\subsection{Coroutine frame, promise object, continuation и безопасност}

При всяка корутина компилаторът създава coroutine frame, в който се съхраняват локалните променливи, текущото състояние на изпълнението и promise object-ът, чрез който корутината комуникира резултат, грешка или завършване. Това има директно отражение върху безопасността: lifetime на frame-а, ownership на резултата и мястото на continuation-а не са второстепенни детайли, а част от correctness contract-а на асинхронния слой.

Трябва да се подчертае, че coroutine frame-ът не е обикновен stack frame. Той може да преживее излизането от текущия стеков контекст и да бъде възобновен по-късно от друга нишка или от event loop инфраструктура. Поради това проблеми като преждевременно разрушаване на frame-а, двойно възобновяване, неправилно предаване на exception състояние или загуба на continuation не са „ниско ниво implementation details“, а фундаментални correctness рискове.

Именно promise object-ът организира връщането на стойност, propagation-а на грешка и поведението при \code{initial\_suspend} и \code{final\_suspend}. От гледна точка на архитектурата това е важен пример за compile-time binding: видът на резултата и semantics на завършването не се определят от динамична registry схема, а са закодирани в coroutine return type-а. Това позволява \code{Task<T>} да бъде проектиран като формална абстракция, а не като неформален „future-подобен“ обект.

\subsection{Корутини, thread pool и event loop инфраструктура}

\code{ThreadPool} и корутините решават различни проблеми. \code{ThreadPool}-ът предоставя ограничен брой worker нишки и policy за планиране на работа, докато корутините описват как една логическа задача може да бъде спряна и по-късно възобновена. Комбинацията между двете позволява блокиращи backend-и да бъдат интегрирани по универсален начин чрез offload, без потребителят да се връща към callback-oriented програмиране.

От друга страна, при истински неблокиращи драйвери корутината може да бъде възобновявана не от worker нишка, а от event loop или callback bridge. Reactor-подобни среди наблюдават готовността на file descriptor-и чрез системни механизми като \code{kqueue} на macOS и BSD системите~\cite{kqueue_man}. Следователно корутината сама по себе си не „изпълнява I/O“, а предоставя формален интерфейс, през който I/O readiness или driver callback може да се превърне в resume действие.

Точно тук става ясно защо корутините не трябва да бъдат смесвани с нишките. Нишката е носител на физическо изпълнение; корутината е носител на логическа последователност. Един и същ асинхронен модел може да използва thread pool в един backend сценарий и event loop в друг, без да променя потребителския синтаксис на \code{co_await}.

\subsection{Значение за Compile-time Object-Relational Mapping библиотеката}

За \ORMName{} този теоретичен анализ има пряко архитектурно следствие. Compile-time моделираният query IR и connector contract определят какво е позволено да се изпълни; coroutine-базираният слой определя как това позволено изчисление може да бъде изпълнено асинхронно. С други думи, типово коректната заявка и коректният execution path са различни, но взаимно зависими пластове.

Това разграничение е особено важно при multi-backend библиотека. PostgreSQL и MySQL предоставят native неблокиращи клиентски интерфейси, но с различни operational форми~\cite{libpq_async,mysql_capi_async}. \code{hiredis} използва callback-базиран async слой~\cite{hiredis}, а драйверът за Cassandra работи чрез future/callback механизъм~\cite{cassandra_futures}. При MongoDB фокусът остава върху безопасен pooling и многонитова употреба на \code{libmongoc}~\cite{mongoc_client_pool}, докато \code{libneo4j-client} следва преобладаващо синхронен request/response модел~\cite{libneo4j_client}. Именно тази нееднородност оправдава асинхронна архитектура с две стратегии: универсален offload модел и native async интеграция за драйвери, които я позволяват.

Следователно корутините не правят автоматично всяка операция неблокираща и не отменят ограниченията на даден драйвер. Те дават формален модел за асинхронно описание на изчислението; реалният performance ефект зависи от това дали backend библиотеката има native non-blocking интерфейс, дали работата се offload-ва към thread pool или дали задачата остава синхронна. Именно затова следващите глави разглеждат отделно compile-time ORM архитектурата, асинхронния слой на рамката и execution semantics на конкретните backend-и.
